\chapter*{Závěr}
\addcontentsline{toc}{chapter}{Závěr}

Tato práce navrhla sadu metrik a~iterativní postup pro
systematické navrhování a~vyhodnocování instrukcí pro AI coding
agenty. Na případové studii systému upomínek faktur ověřila
proveditelnost tohoto postupu na jednom projektu a~ablacemi
prozkoumala, které složky instrukcí se podílely na měřeném chování
agenta. Studie zahrnovala pět pilotních iterací a~dvě ablace, každou
ověřenou dvěma nezávislými běhy. Její závěry proto mají indikativní
povahu případové studie, ne podobu statistické generalizace.

Prvním cílem bylo navrhnout sadu metrik pokrývající proces, kvalitu
produktu a~zdroje. Výsledná sada 19~metrik zachytila rozdíly, které
by samotné hodnocení funkční korektnosti spojilo dohromady. Procesní
metriky ukázaly, zda agent dodržoval očekávaný pracovní postup
a~zanechal auditovatelnou stopu. Produktové metriky zachytily
korektnost a~kvalitu vzniklého kódu. Metriky zdrojů doplnily pohled
na cenu, čas a~práci s~kontextem. Kombinace deterministických
a~judge-based metrik se ukázala jako prakticky použitelná, protože
jednotlivé typy měření zachycovaly odlišné vlastnosti běhu.

Druhým cílem bylo ověřit proveditelnost iterativního postupu návrhu
instrukcí řízeného těmito metrikami. Případová studie ukázala, že
metriky lze použít nejen k~vyhodnocení výsledku, ale i~k~diagnóze
selhání a~návrhu další úpravy instrukcí. Mezi běhy se objevila
zlepšení i~regrese, což potvrdilo potřebu opakovaného měření celé
sady metrik. Nejstabilnějším vzorcem úspěšných úprav bylo převedení
obecného pravidla na konkrétní příkaz a~následně na verifikační krok.
Tento vzorec byl nejsilnější tam, kde existoval deterministický
nástroj s~jednoznačnou zpětnou vazbou.

Třetím cílem bylo ablacemi prozkoumat, které složky instrukcí
vytvořených v~cíli~2 přispívají k~měřenému chování agenta a~které
jsou redundantní. Výsledky byly konzistentní s~interpretací, že
verifikační kroky byly v~této případové studii prakticky důležitou
složkou pracovního postupu. Jejich odebrání zhoršilo vybrané metriky
korektnosti i~kvality. Sekce Package Quality se v~deterministických
metrikách jevila převážně redundantní, ale její odebrání se projevilo
v~designové kvalitě hodnocené LLM-as-judge. Ablace zároveň ukázaly,
že dodržení vícekrokového pracovního postupu zůstává citlivé na
nedeterminismus modelu.

Hlavním přínosem práce je propojení sady metrik s~iterativním
postupem návrhu instrukcí. Přenositelný není konkrétní soubor
\texttt{AGENTS.md} ani naměřené hodnoty, ale způsob práce, tedy
zaznamenat běh, vyhodnotit jej ve více dimenzích, diagnostikovat
selhání a~upravit instrukce podle naměřených dat. Konkrétní prahy
metrik a~znění instrukcí je nutné přizpůsobit modelu, nástroji
a~projektu.

Práce tím ukazuje, že instrukce pro AI coding agenty lze zkoumat jako
samostatně měřenou nezávislou proměnnou, nikoli jen jako doplněk
promptu. Pro praktické použití z~toho plyne, že nestačí sledovat,
zda agent úkol vyřešil. Stejně důležité je měřit, jak při tom
postupoval, jak kvalitní kód zanechal a~kolik zdrojů spotřeboval.
